\documentclass[11pt]{article}
\usepackage[margin=1in]{geometry}
\usepackage{amsmath}
\usepackage{amssymb}
\usepackage{booktabs}
\usepackage{graphicx}
\usepackage{longtable}
\usepackage[table]{xcolor}
\usepackage[colorlinks=true,urlcolor=blue,linkcolor=black]{hyperref}
\setlength{\parskip}{0.5em}
\setlength{\parindent}{0pt}
% \texttt{ring\_count} and friends are unbreakable and otherwise poke into the
% margin; let TeX stretch interword space instead.
\sloppy

\title{The $(\lambda, N)$ sweep on ring count\\\large \texttt{marginal} against \texttt{edlm}, MDLM and UDLM}
\author{}\date{2026-08-15}
\begin{document}\maketitle
\section{Setup}

252 runs, 1{,}024 samples each at $T=32$ steps, seed 1, reward
\texttt{ring\_count}. Grid: $\lambda \in \{0, 0.5, 1, 2, 5, 10, 30, 100, 1000\}$ against
$N \in \{10, 30, 100, 300, 500, 1000, 2000\}$, for both mixture-sampling modes
and both base models.

$N$ is the number of $\mathbf{x}_0$ candidates drawn per denoising step --- the
Monte Carlo budget of the estimator. $\lambda$ is the tilt in
$q(\mathbf{x}_0) \propto p_\theta(\mathbf{x}_0)\,e^{\lambda r(\mathbf{x}_0)}$.
The two modes differ only in how $\mathbf{x}_{t-1}$ is drawn from the resulting
mixture: \texttt{marginal} averages the $N$ individually normalised posteriors,
\texttt{edlm} draws a component and then samples from it. They share
per-position marginals and differ only in the joint.

\section{Reading ESS: use the median, never the last step}

Effective sample size $\mathrm{ESS} = 1/\sum_n w_n^2$ says which end of the grid
a run is on: $\mathrm{ESS}=N$ means the tilt is doing nothing, $\mathrm{ESS}\to1$
means it has collapsed onto a single candidate.

The driver script prints only the \emph{final} step's ESS, and that number is
actively misleading on MDLM. Absorbing-state sampling has almost everything
unmasked by the last step, so all $N$ candidates decode to the same sequence and
tie on reward --- the last-step ESS reads $\approx N$ however hard $\lambda$ is
tilting. Judged on it, MDLM looks entirely unguided. Its novel QED in fact climbs
across the $\lambda$ grid. Every ESS below is the \textbf{median over the
trajectory}, digested from the per-run logs into
\texttt{results/\_ess\_summary.csv}.

\begin{table}[htbp]
\centering
\small
\resizebox{\ifdim\width>\linewidth\linewidth\else\width\fi}{!}{%
\begin{tabular}{rrrrrrrrrr}
\toprule
$N$ & $\lambda=0$ & $\lambda=0.5$ & $\lambda=1$ & $\lambda=2$ & $\lambda=5$ & $\lambda=10$ & $\lambda=30$ & $\lambda=100$ & $\lambda=1000$ \\
\midrule
10 & 1.00 & 0.97 & 0.94 & 0.90 & 0.88 & 0.88 & 0.88 & 0.88 & 0.88 \\
30 & 1.00 & 0.96 & 0.91 & 0.84 & 0.80 & 0.81 & 0.80 & 0.80 & 0.80 \\
100 & 1.00 & 0.96 & 0.89 & 0.80 & 0.73 & 0.71 & 0.71 & 0.71 & 0.71 \\
300 & 1.00 & 0.95 & 0.86 & 0.74 & 0.65 & 0.65 & 0.64 & 0.65 & 0.65 \\
500 & 1.00 & 0.95 & 0.85 & 0.72 & 0.63 & 0.62 & 0.62 & 0.62 & 0.62 \\
1000 & 0.00 & 0.94 & 0.83 & 0.70 & 0.59 & 0.58 & 0.58 & 0.58 & 0.58 \\
2000 & 0.00 & 0.92 & 0.78 & 0.61 & 0.51 & 0.49 & 0.49 & 0.49 & 0.49 \\
\bottomrule
\end{tabular}}
\caption{MDLM, \texttt{marginal}: median ESS as a fraction of $N$.}
\label{tab:ess-ring_count-MDLM-marginal}
\end{table}

\begin{table}[htbp]
\centering
\small
\resizebox{\ifdim\width>\linewidth\linewidth\else\width\fi}{!}{%
\begin{tabular}{rrrrrrrrrr}
\toprule
$N$ & $\lambda=0$ & $\lambda=0.5$ & $\lambda=1$ & $\lambda=2$ & $\lambda=5$ & $\lambda=10$ & $\lambda=30$ & $\lambda=100$ & $\lambda=1000$ \\
\midrule
10 & 1.00 & 0.97 & 0.94 & 0.91 & 0.88 & 0.88 & 0.88 & 0.88 & 0.88 \\
30 & 1.00 & 0.96 & 0.91 & 0.86 & 0.81 & 0.80 & 0.80 & 0.80 & 0.80 \\
100 & 1.00 & 0.96 & 0.89 & 0.80 & 0.72 & 0.70 & 0.69 & 0.69 & 0.69 \\
300 & 1.00 & 0.95 & 0.87 & 0.75 & 0.64 & 0.65 & 0.64 & 0.64 & 0.64 \\
500 & 1.00 & 0.95 & 0.86 & 0.72 & 0.64 & 0.63 & 0.63 & 0.63 & 0.63 \\
1000 & 0.00 & 0.94 & 0.82 & 0.68 & 0.60 & 0.59 & 0.58 & 0.58 & 0.58 \\
2000 & 0.00 & 0.93 & 0.79 & 0.62 & 0.53 & 0.51 & 0.50 & 0.50 & 0.50 \\
\bottomrule
\end{tabular}}
\caption{MDLM, \texttt{edlm}: median ESS as a fraction of $N$.}
\label{tab:ess-ring_count-MDLM-edlm}
\end{table}

\begin{table}[htbp]
\centering
\small
\resizebox{\ifdim\width>\linewidth\linewidth\else\width\fi}{!}{%
\begin{tabular}{rrrrrrrrrr}
\toprule
$N$ & $\lambda=0$ & $\lambda=0.5$ & $\lambda=1$ & $\lambda=2$ & $\lambda=5$ & $\lambda=10$ & $\lambda=30$ & $\lambda=100$ & $\lambda=1000$ \\
\midrule
10 & 1.00 & 0.98 & 0.94 & 0.88 & 0.85 & 0.84 & 0.85 & 0.84 & 0.84 \\
30 & 1.00 & 0.97 & 0.91 & 0.80 & 0.72 & 0.72 & 0.72 & 0.72 & 0.72 \\
100 & 1.00 & 0.97 & 0.88 & 0.71 & 0.51 & 0.51 & 0.50 & 0.50 & 0.50 \\
300 & 1.00 & 0.96 & 0.86 & 0.57 & 0.26 & 0.26 & 0.26 & 0.26 & 0.26 \\
500 & 1.00 & 0.97 & 0.85 & 0.52 & 0.19 & 0.19 & 0.18 & 0.18 & 0.18 \\
1000 & 0.00 & 0.97 & 0.85 & 0.43 & 0.11 & 0.11 & 0.11 & 0.11 & 0.11 \\
2000 & 0.00 & 0.97 & 0.82 & 0.36 & 0.07 & 0.07 & 0.07 & 0.07 & 0.07 \\
\bottomrule
\end{tabular}}
\caption{UDLM, \texttt{marginal}: median ESS as a fraction of $N$.}
\label{tab:ess-ring_count-UDLM-marginal}
\end{table}

\begin{table}[htbp]
\centering
\small
\resizebox{\ifdim\width>\linewidth\linewidth\else\width\fi}{!}{%
\begin{tabular}{rrrrrrrrrr}
\toprule
$N$ & $\lambda=0$ & $\lambda=0.5$ & $\lambda=1$ & $\lambda=2$ & $\lambda=5$ & $\lambda=10$ & $\lambda=30$ & $\lambda=100$ & $\lambda=1000$ \\
\midrule
10 & 1.00 & 0.98 & 0.94 & 0.88 & 0.85 & 0.85 & 0.85 & 0.85 & 0.85 \\
30 & 1.00 & 0.97 & 0.91 & 0.81 & 0.73 & 0.73 & 0.73 & 0.73 & 0.73 \\
100 & 1.00 & 0.97 & 0.89 & 0.71 & 0.51 & 0.51 & 0.51 & 0.51 & 0.51 \\
300 & 1.00 & 0.96 & 0.85 & 0.58 & 0.27 & 0.27 & 0.27 & 0.27 & 0.27 \\
500 & 1.00 & 0.97 & 0.85 & 0.51 & 0.19 & 0.18 & 0.18 & 0.18 & 0.18 \\
1000 & 0.00 & 0.96 & 0.83 & 0.42 & 0.11 & 0.11 & 0.11 & 0.11 & 0.11 \\
2000 & 0.00 & 0.96 & 0.83 & 0.36 & 0.07 & 0.07 & 0.07 & 0.07 & 0.07 \\
\bottomrule
\end{tabular}}
\caption{UDLM, \texttt{edlm}: median ESS as a fraction of $N$.}
\label{tab:ess-ring_count-UDLM-edlm}
\end{table}

\section{The grid}
\begin{table}[htbp]
\centering
\small
\resizebox{\ifdim\width>\linewidth\linewidth\else\width\fi}{!}{%
\begin{tabular}{rrrrrrrrrr}
\toprule
$N$ & $\lambda=0$ & $\lambda=0.5$ & $\lambda=1$ & $\lambda=2$ & $\lambda=5$ & $\lambda=10$ & $\lambda=30$ & $\lambda=100$ & $\lambda=1000$ \\
\midrule
10 & 2.477 & 2.751 & 2.824 & 2.815 & 2.819 & 2.820 & 2.820 & 2.820 & 2.820 \\
30 & 2.365 & 2.552 & 2.657 & 2.683 & 2.681 & 2.639 & 2.633 & 2.633 & 2.633 \\
100 & 2.346 & 2.573 & 2.822 & 2.979 & 3.080 & 3.129 & 3.108 & 3.108 & 3.108 \\
300 & 2.468 & 2.716 & 3.033 & 3.260 & 3.345 & 3.351 & 3.367 & 3.367 & 3.367 \\
500 & 2.548 & 2.698 & 2.952 & 3.173 & 3.311 & 3.413 & 3.410 & 3.410 & 3.410 \\
1000 & 2.426 & 2.603 & 2.911 & 3.310 & 3.452 & 3.499 & 3.493 & 3.493 & 3.493 \\
2000 & 2.379 & 2.672 & 3.045 & 3.615 & 3.715 & 3.611 & 3.627 & 3.627 & 3.627 \\
\bottomrule
\end{tabular}}
\caption{MDLM, \texttt{marginal}: novel-ring count mean.}
\label{tab:grid-ring_count-MDLM-marginal}
\end{table}

\begin{table}[htbp]
\centering
\small
\resizebox{\ifdim\width>\linewidth\linewidth\else\width\fi}{!}{%
\begin{tabular}{rrrrrrrrrr}
\toprule
$N$ & $\lambda=0$ & $\lambda=0.5$ & $\lambda=1$ & $\lambda=2$ & $\lambda=5$ & $\lambda=10$ & $\lambda=30$ & $\lambda=100$ & $\lambda=1000$ \\
\midrule
10 & 2.567 & 2.707 & 2.810 & 2.864 & 2.783 & 2.773 & 2.770 & 2.770 & 2.770 \\
30 & 2.697 & 2.841 & 2.874 & 2.899 & 2.797 & 2.797 & 2.797 & 2.797 & 2.797 \\
100 & 2.284 & 2.550 & 2.854 & 2.834 & 3.039 & 3.107 & 3.103 & 3.103 & 3.103 \\
300 & 2.493 & 2.736 & 2.964 & 3.080 & 3.396 & 3.396 & 3.396 & 3.396 & 3.396 \\
500 & 2.318 & 2.724 & 3.107 & 3.146 & 3.243 & 3.194 & 3.184 & 3.184 & 3.184 \\
1000 & 2.612 & 2.803 & 3.174 & 3.306 & 3.484 & 3.543 & 3.538 & 3.538 & 3.538 \\
2000 & 2.718 & 2.840 & 3.090 & 3.546 & 3.793 & 3.754 & 3.752 & 3.752 & 3.752 \\
\bottomrule
\end{tabular}}
\caption{MDLM, \texttt{edlm}: novel-ring count mean.}
\label{tab:grid-ring_count-MDLM-edlm}
\end{table}

\begin{table}[htbp]
\centering
\small
\resizebox{\ifdim\width>\linewidth\linewidth\else\width\fi}{!}{%
\begin{tabular}{rrrrrrrrrr}
\toprule
$N$ & $\lambda=0$ & $\lambda=0.5$ & $\lambda=1$ & $\lambda=2$ & $\lambda=5$ & $\lambda=10$ & $\lambda=30$ & $\lambda=100$ & $\lambda=1000$ \\
\midrule
10 & 1.656 & 1.865 & 1.935 & 1.873 & 1.868 & 1.884 & 1.884 & 1.884 & 1.884 \\
30 & 1.708 & 1.785 & 1.787 & 1.884 & 1.789 & 1.799 & 1.804 & 1.804 & 1.804 \\
100 & 1.660 & 1.776 & 1.825 & 1.846 & 1.920 & 1.936 & 1.932 & 1.932 & 1.932 \\
300 & 1.637 & 1.768 & 1.764 & 1.954 & 2.168 & 2.110 & 2.125 & 2.125 & 2.125 \\
500 & 1.592 & 1.627 & 1.721 & 2.008 & 2.317 & 2.162 & 2.211 & 2.211 & 2.211 \\
1000 & 1.543 & 1.646 & 1.842 & 2.129 & 2.302 & 2.364 & 2.351 & 2.351 & 2.351 \\
2000 & 1.469 & 1.654 & 1.733 & 1.953 & 2.524 & 2.572 & 2.581 & 2.581 & 2.581 \\
\bottomrule
\end{tabular}}
\caption{UDLM, \texttt{marginal}: novel-ring count mean.}
\label{tab:grid-ring_count-UDLM-marginal}
\end{table}

\begin{table}[htbp]
\centering
\small
\resizebox{\ifdim\width>\linewidth\linewidth\else\width\fi}{!}{%
\begin{tabular}{rrrrrrrrrr}
\toprule
$N$ & $\lambda=0$ & $\lambda=0.5$ & $\lambda=1$ & $\lambda=2$ & $\lambda=5$ & $\lambda=10$ & $\lambda=30$ & $\lambda=100$ & $\lambda=1000$ \\
\midrule
10 & 1.624 & 1.764 & 1.775 & 1.828 & 1.738 & 1.732 & 1.732 & 1.732 & 1.732 \\
30 & 1.524 & 1.707 & 1.853 & 1.718 & 1.734 & 1.717 & 1.717 & 1.717 & 1.717 \\
100 & 1.600 & 1.801 & 1.888 & 1.860 & 1.918 & 1.791 & 1.793 & 1.793 & 1.793 \\
300 & 1.693 & 1.801 & 1.820 & 2.028 & 2.161 & 2.043 & 2.027 & 2.027 & 2.027 \\
500 & 1.653 & 1.901 & 1.902 & 1.933 & 2.195 & 2.199 & 2.180 & 2.180 & 2.180 \\
1000 & 1.799 & 1.888 & 1.966 & 2.328 & 2.492 & 2.401 & 2.392 & 2.392 & 2.392 \\
2000 & 1.797 & 2.013 & 1.945 & 2.077 & 2.540 & 2.717 & 2.727 & 2.727 & 2.727 \\
\bottomrule
\end{tabular}}
\caption{UDLM, \texttt{edlm}: novel-ring count mean.}
\label{tab:grid-ring_count-UDLM-edlm}
\end{table}

\begin{table}[htbp]
\centering
\small
\resizebox{\ifdim\width>\linewidth\linewidth\else\width\fi}{!}{%
\begin{tabular}{rrrrrrrrrr}
\toprule
$N$ & $\lambda=0$ & $\lambda=0.5$ & $\lambda=1$ & $\lambda=2$ & $\lambda=5$ & $\lambda=10$ & $\lambda=30$ & $\lambda=100$ & $\lambda=1000$ \\
\midrule
10 & 216 & 233 & 244 & 243 & 254 & 256 & 256 & 256 & 256 \\
30 & 211 & 232 & 245 & 265 & 276 & 288 & 289 & 289 & 289 \\
100 & 211 & 232 & 259 & 281 & 287 & 279 & 279 & 279 & 279 \\
300 & 220 & 236 & 241 & 281 & 293 & 316 & 316 & 316 & 316 \\
500 & 228 & 245 & 272 & 294 & 331 & 334 & 334 & 334 & 334 \\
1000 & 204 & 229 & 257 & 281 & 361 & 367 & 369 & 369 & 369 \\
2000 & 190 & 204 & 245 & 304 & 354 & 357 & 362 & 362 & 362 \\
\bottomrule
\end{tabular}}
\caption{MDLM, \texttt{marginal}: number of novel molecules (out of 1{,}024 generated).}
\label{tab:novel-ring_count-MDLM-marginal}
\end{table}

\begin{table}[htbp]
\centering
\small
\resizebox{\ifdim\width>\linewidth\linewidth\else\width\fi}{!}{%
\begin{tabular}{rrrrrrrrrr}
\toprule
$N$ & $\lambda=0$ & $\lambda=0.5$ & $\lambda=1$ & $\lambda=2$ & $\lambda=5$ & $\lambda=10$ & $\lambda=30$ & $\lambda=100$ & $\lambda=1000$ \\
\midrule
10 & 215 & 242 & 247 & 243 & 254 & 256 & 256 & 256 & 256 \\
30 & 234 & 258 & 262 & 258 & 256 & 251 & 251 & 251 & 251 \\
100 & 208 & 242 & 267 & 289 & 307 & 309 & 310 & 310 & 310 \\
300 & 209 & 231 & 247 & 275 & 285 & 288 & 288 & 288 & 288 \\
500 & 195 & 214 & 252 & 294 & 300 & 319 & 320 & 320 & 320 \\
1000 & 219 & 249 & 281 & 307 & 339 & 363 & 364 & 364 & 364 \\
2000 & 213 & 231 & 266 & 328 & 367 & 378 & 379 & 379 & 379 \\
\bottomrule
\end{tabular}}
\caption{MDLM, \texttt{edlm}: number of novel molecules (out of 1{,}024 generated).}
\label{tab:novel-ring_count-MDLM-edlm}
\end{table}

\begin{table}[htbp]
\centering
\small
\resizebox{\ifdim\width>\linewidth\linewidth\else\width\fi}{!}{%
\begin{tabular}{rrrrrrrrrr}
\toprule
$N$ & $\lambda=0$ & $\lambda=0.5$ & $\lambda=1$ & $\lambda=2$ & $\lambda=5$ & $\lambda=10$ & $\lambda=30$ & $\lambda=100$ & $\lambda=1000$ \\
\midrule
10 & 189 & 208 & 230 & 251 & 243 & 241 & 241 & 241 & 241 \\
30 & 195 & 214 & 202 & 242 & 251 & 264 & 265 & 265 & 265 \\
100 & 194 & 214 & 229 & 240 & 286 & 267 & 265 & 265 & 265 \\
300 & 193 & 207 & 212 & 240 & 268 & 292 & 288 & 288 & 288 \\
500 & 196 & 217 & 222 & 249 & 278 & 284 & 279 & 279 & 279 \\
1000 & 186 & 206 & 228 & 241 & 288 & 294 & 308 & 308 & 308 \\
2000 & 207 & 234 & 243 & 256 & 275 & 278 & 277 & 277 & 277 \\
\bottomrule
\end{tabular}}
\caption{UDLM, \texttt{marginal}: number of novel molecules (out of 1{,}024 generated).}
\label{tab:novel-ring_count-UDLM-marginal}
\end{table}

\begin{table}[htbp]
\centering
\small
\resizebox{\ifdim\width>\linewidth\linewidth\else\width\fi}{!}{%
\begin{tabular}{rrrrrrrrrr}
\toprule
$N$ & $\lambda=0$ & $\lambda=0.5$ & $\lambda=1$ & $\lambda=2$ & $\lambda=5$ & $\lambda=10$ & $\lambda=30$ & $\lambda=100$ & $\lambda=1000$ \\
\midrule
10 & 165 & 182 & 204 & 233 & 252 & 257 & 257 & 257 & 257 \\
30 & 191 & 222 & 238 & 266 & 274 & 276 & 276 & 276 & 276 \\
100 & 180 & 211 & 233 & 265 & 269 & 296 & 295 & 295 & 295 \\
300 & 179 & 191 & 205 & 251 & 310 & 301 & 301 & 301 & 301 \\
500 & 190 & 222 & 214 & 252 & 277 & 297 & 295 & 295 & 295 \\
1000 & 164 & 215 & 235 & 271 & 301 & 309 & 309 & 309 & 309 \\
2000 & 212 & 239 & 253 & 271 & 313 & 311 & 319 & 319 & 319 \\
\bottomrule
\end{tabular}}
\caption{UDLM, \texttt{edlm}: number of novel molecules (out of 1{,}024 generated).}
\label{tab:novel-ring_count-UDLM-edlm}
\end{table}

\begin{table}[htbp]
\centering
\small
\resizebox{\ifdim\width>\linewidth\linewidth\else\width\fi}{!}{%
\begin{tabular}{rrrrrrrrrr}
\toprule
$N$ & $\lambda=0$ & $\lambda=0.5$ & $\lambda=1$ & $\lambda=2$ & $\lambda=5$ & $\lambda=10$ & $\lambda=30$ & $\lambda=100$ & $\lambda=1000$ \\
\midrule
10 & 9.000 & 9.000 & 10.000 & 10.000 & 10.000 & 10.000 & 10.000 & 10.000 & 10.000 \\
30 & 8.000 & 10.000 & 8.000 & 8.000 & 8.000 & 8.000 & 8.000 & 8.000 & 8.000 \\
100 & 9.000 & 10.000 & 9.000 & 10.000 & 10.000 & 10.000 & 10.000 & 10.000 & 10.000 \\
300 & 7.000 & 14.000 & 14.000 & 14.000 & 14.000 & 14.000 & 14.000 & 14.000 & 14.000 \\
500 & 13.000 & 13.000 & 13.000 & 13.000 & 13.000 & 13.000 & 13.000 & 13.000 & 13.000 \\
1000 & 9.000 & 9.000 & 9.000 & 10.000 & 10.000 & 10.000 & 10.000 & 10.000 & 10.000 \\
2000 & 8.000 & 10.000 & 10.000 & 12.000 & 12.000 & 12.000 & 12.000 & 12.000 & 12.000 \\
\bottomrule
\end{tabular}}
\\[0.4em] \begin{minipage}{0.92\textwidth}\footnotesize A raw max grows with the number of molecules it is taken over, and these cells differ several-fold in novel count, so a higher max here is partly a deeper pool. That is a real advantage at a fixed generation budget -- every cell generated the same 1,024 sequences -- but it is not a claim about distribution shape. Section~\ref{sec:beyond-mean-ring_count} separates the two.\end{minipage}
\caption{MDLM, \texttt{marginal}: highest novel-ring count of the run.}
\label{tab:max-ring_count-MDLM-marginal}
\end{table}

\begin{table}[htbp]
\centering
\small
\resizebox{\ifdim\width>\linewidth\linewidth\else\width\fi}{!}{%
\begin{tabular}{rrrrrrrrrr}
\toprule
$N$ & $\lambda=0$ & $\lambda=0.5$ & $\lambda=1$ & $\lambda=2$ & $\lambda=5$ & $\lambda=10$ & $\lambda=30$ & $\lambda=100$ & $\lambda=1000$ \\
\midrule
10 & 7.000 & 9.000 & 9.000 & 9.000 & 9.000 & 9.000 & 9.000 & 9.000 & 9.000 \\
30 & 9.000 & 10.000 & 10.000 & 10.000 & 10.000 & 10.000 & 10.000 & 10.000 & 10.000 \\
100 & 8.000 & 10.000 & 10.000 & 10.000 & 10.000 & 10.000 & 10.000 & 10.000 & 10.000 \\
300 & 8.000 & 10.000 & 10.000 & 10.000 & 9.000 & 9.000 & 9.000 & 9.000 & 9.000 \\
500 & 10.000 & 10.000 & 10.000 & 10.000 & 10.000 & 10.000 & 10.000 & 10.000 & 10.000 \\
1000 & 7.000 & 9.000 & 10.000 & 10.000 & 11.000 & 11.000 & 11.000 & 11.000 & 11.000 \\
2000 & 8.000 & 10.000 & 10.000 & 10.000 & 10.000 & 10.000 & 10.000 & 10.000 & 10.000 \\
\bottomrule
\end{tabular}}
\caption{MDLM, \texttt{edlm}: highest novel-ring count of the run.}
\label{tab:max-ring_count-MDLM-edlm}
\end{table}

\begin{table}[htbp]
\centering
\small
\resizebox{\ifdim\width>\linewidth\linewidth\else\width\fi}{!}{%
\begin{tabular}{rrrrrrrrrr}
\toprule
$N$ & $\lambda=0$ & $\lambda=0.5$ & $\lambda=1$ & $\lambda=2$ & $\lambda=5$ & $\lambda=10$ & $\lambda=30$ & $\lambda=100$ & $\lambda=1000$ \\
\midrule
10 & 6.000 & 7.000 & 7.000 & 7.000 & 7.000 & 7.000 & 7.000 & 7.000 & 7.000 \\
30 & 6.000 & 7.000 & 7.000 & 7.000 & 6.000 & 6.000 & 6.000 & 6.000 & 6.000 \\
100 & 7.000 & 6.000 & 6.000 & 6.000 & 7.000 & 8.000 & 8.000 & 8.000 & 8.000 \\
300 & 5.000 & 7.000 & 8.000 & 7.000 & 7.000 & 6.000 & 6.000 & 6.000 & 6.000 \\
500 & 6.000 & 6.000 & 5.000 & 7.000 & 7.000 & 7.000 & 7.000 & 7.000 & 7.000 \\
1000 & 5.000 & 10.000 & 10.000 & 7.000 & 10.000 & 7.000 & 7.000 & 7.000 & 7.000 \\
2000 & 5.000 & 6.000 & 5.000 & 7.000 & 6.000 & 10.000 & 10.000 & 10.000 & 10.000 \\
\bottomrule
\end{tabular}}
\caption{UDLM, \texttt{marginal}: highest novel-ring count of the run.}
\label{tab:max-ring_count-UDLM-marginal}
\end{table}

\begin{table}[htbp]
\centering
\small
\resizebox{\ifdim\width>\linewidth\linewidth\else\width\fi}{!}{%
\begin{tabular}{rrrrrrrrrr}
\toprule
$N$ & $\lambda=0$ & $\lambda=0.5$ & $\lambda=1$ & $\lambda=2$ & $\lambda=5$ & $\lambda=10$ & $\lambda=30$ & $\lambda=100$ & $\lambda=1000$ \\
\midrule
10 & 6.000 & 8.000 & 6.000 & 6.000 & 5.000 & 5.000 & 5.000 & 5.000 & 5.000 \\
30 & 7.000 & 7.000 & 7.000 & 7.000 & 7.000 & 7.000 & 7.000 & 7.000 & 7.000 \\
100 & 6.000 & 6.000 & 7.000 & 6.000 & 5.000 & 6.000 & 6.000 & 6.000 & 6.000 \\
300 & 7.000 & 7.000 & 7.000 & 8.000 & 8.000 & 7.000 & 7.000 & 7.000 & 7.000 \\
500 & 7.000 & 9.000 & 7.000 & 7.000 & 10.000 & 7.000 & 7.000 & 7.000 & 7.000 \\
1000 & 6.000 & 6.000 & 6.000 & 7.000 & 8.000 & 8.000 & 8.000 & 8.000 & 8.000 \\
2000 & 7.000 & 7.000 & 6.000 & 7.000 & 10.000 & 10.000 & 10.000 & 10.000 & 10.000 \\
\bottomrule
\end{tabular}}
\caption{UDLM, \texttt{edlm}: highest novel-ring count of the run.}
\label{tab:max-ring_count-UDLM-edlm}
\end{table}

\section{\texttt{marginal} against \texttt{edlm}}
\begin{table}[htbp]
\centering
\small
\resizebox{\ifdim\width>\linewidth\linewidth\else\width\fi}{!}{%
\begin{tabular}{lr rrrr}
\toprule
model & $\lambda$ & mean $|\Delta|$ property & mean $\Delta$ property & mean $|\Delta|$ novel & mean $\Delta$ novel \\
\midrule
MDLM & 0 & 0.181 & -0.097 & 15.6 & -1.9 \\
MDLM & 0.5 & 0.110 & -0.091 & 18.3 & -8.0 \\
MDLM & 1 & 0.114 & -0.090 & 14.1 & -8.4 \\
MDLM & 2 & 0.099 & +0.023 & 10.1 & -6.4 \\
MDLM & 5 & 0.060 & -0.019 & 16.3 & +6.9 \\
MDLM & 10 & 0.097 & -0.015 & 19.3 & +4.7 \\
MDLM & 30 & 0.092 & -0.012 & 19.0 & +5.3 \\
MDLM & 100 & 0.092 & -0.012 & 19.0 & +5.3 \\
MDLM & 1000 & 0.092 & -0.012 & 19.0 & +5.3 \\
UDLM & 0 & 0.139 & -0.061 & 12.7 & +11.3 \\
UDLM & 0.5 & 0.159 & -0.108 & 10.3 & +2.6 \\
UDLM & 1 & 0.123 & -0.077 & 14.0 & -2.3 \\
UDLM & 2 & 0.100 & -0.018 & 18.0 & -12.9 \\
UDLM & 5 & 0.074 & +0.016 & 20.4 & -15.3 \\
UDLM & 10 & 0.095 & +0.032 & 18.1 & -18.1 \\
UDLM & 30 & 0.099 & +0.046 & 18.4 & -18.4 \\
UDLM & 100 & 0.099 & +0.046 & 18.4 & -18.4 \\
UDLM & 1000 & 0.099 & +0.046 & 18.4 & -18.4 \\
\bottomrule
\end{tabular}}
\\[0.4em] \begin{minipage}{0.92\textwidth}\footnotesize A systematic winner would show as a consistent sign in the $\Delta$ columns. There is none.\end{minipage}
\caption{\texttt{marginal} minus \texttt{edlm}, averaged over the five $N$ at each $\lambda$. $|\Delta|$ is the typical size of the gap, $\Delta$ its sign.}
\label{tab:modes-ring_count}
\end{table}


Two things were predicted before the sweep ran, and only one holds.

\textbf{Predicted and confirmed:} the two modes are empirically
interchangeable. Across the whole grid the mean absolute gap in novel-ring count is
0.107 with mean signed gap
-0.022 --- no consistent winner, at any $\lambda$ or $N$.

\textbf{Predicted and \emph{not} confirmed:} the gap was expected to
\emph{shrink} as $\lambda$ grows, because one-hot weights make
$\sum_n w_n q(\cdot\,|\,\mathbf{x}_0^{(n)})$ collapse to the single component
\texttt{edlm} would have drawn, and a posterior conditioned on a fixed
$\mathbf{x}_0$ already factorises over positions. Measured, the gap is flat:
0.160 at $\lambda=0$ against
0.096 at $\lambda=1000$. The reason is
visible in the ESS tables --- the weights never actually become one-hot in this
grid, so the limit is never entered. The flat gap is therefore sampling noise
between two samplers that agree, which is the same conclusion by a weaker route.

\section{$\lambda = 0$: the unguided control}

At $\lambda=0$ the weights are uniform, so no tilt is applied whatever $N$ or the
mode is. Every cell of Table~\ref{tab:lam0-ring_count} should therefore be
the base sampler, and the spread across it measures the sampling noise that every
other comparison in this report has to clear.

\begin{table}[htbp]
\centering
\small
\resizebox{\ifdim\width>\linewidth\linewidth\else\width\fi}{!}{%
\begin{tabular}{llrrrrrrr}
\toprule
model & mode & $N=10$ & $N=30$ & $N=100$ & $N=300$ & $N=500$ & $N=1000$ & $N=2000$ \\
\midrule
MDLM & marginal & 2.477 & 2.365 & 2.346 & 2.468 & 2.548 & 2.426 & 2.379 \\
MDLM & edlm & 2.567 & 2.697 & 2.284 & 2.493 & 2.318 & 2.612 & 2.718 \\
UDLM & marginal & 1.656 & 1.708 & 1.660 & 1.637 & 1.592 & 1.543 & 1.469 \\
UDLM & edlm & 1.624 & 1.524 & 1.600 & 1.693 & 1.653 & 1.799 & 1.797 \\
\bottomrule
\end{tabular}}
\\[0.4em] \begin{minipage}{0.92\textwidth}\footnotesize This is the implementation control: at $\lambda=0$ the weights are uniform, so no tilt is applied and the numbers must be flat in $N$ up to sampling noise. They are.\end{minipage}
\caption{novel-ring count at $\lambda=0$. Every entry should be the base sampler, independent of $N$ and of the mode.}
\label{tab:lam0-ring_count}
\end{table}


This table also calibrates everything else in the report. The widest $\lambda=0$
spread within a single model is 0.435 in novel-ring count --- runs that are
identical in distribution, differing only by sampling noise. The mean gap between
\texttt{marginal} and \texttt{edlm} across the whole grid is 0.107,
\textbf{smaller} than that floor. So the two
modes are not merely close: their difference is
indistinguishable from noise.

\section{The frontier against D-CBG}

Both methods are swept, then compared on the axes the paper's Figure 3 uses:
number of novel molecules against novel-ring count mean. Comparing at single
hyperparameter points would be meaningless --- the paper tunes $\gamma$ per model
\emph{and} per property.

\begin{table}[htbp]
\centering
\small
\resizebox{\ifdim\width>\linewidth\linewidth\else\width\fi}{!}{%
\begin{tabular}{ll rrr}
\toprule
method & setting & valid & novel & novel-ring count \\
\midrule
ours & edlm, $N=2000$, $\lambda=5$ & 766 & 367 & 3.793 \\
ours & edlm, $N=2000$, $\lambda=10$ & 754 & 378 & 3.754 \\
ours & edlm, $N=2000$, $\lambda=30$ & 754 & 379 & 3.752 \\
ours & edlm, $N=2000$, $\lambda=100$ & 754 & 379 & 3.752 \\
ours & edlm, $N=2000$, $\lambda=1000$ & 754 & 379 & 3.752 \\
D-CBG & $\gamma=1$, exact & 188 & 149 & 4.745 \\
D-CBG & $\gamma=3$, approx. & 267 & 150 & 3.400 \\
D-CBG & $\gamma=2$, approx. & 352 & 176 & 2.966 \\
D-CBG & $\gamma=1$, approx. & 493 & 209 & 2.809 \\
\bottomrule
\end{tabular}}
\\[0.4em] \begin{minipage}{0.92\textwidth}\footnotesize 1 collapsed run(s) with fewer than 10 novel molecules were excluded; a property mean over a handful of molecules is not an operating point.\end{minipage}
\caption{MDLM: the Pareto-optimal settings of each method on (novel count, novel-ring count).}
\label{tab:frontier-ring_count-MDLM}
\end{table}

\begin{table}[htbp]
\centering
\small
\resizebox{\ifdim\width>\linewidth\linewidth\else\width\fi}{!}{%
\begin{tabular}{ll rrr}
\toprule
method & setting & valid & novel & novel-ring count \\
\midrule
ours & edlm, $N=2000$, $\lambda=30$ & 982 & 319 & 2.727 \\
ours & edlm, $N=2000$, $\lambda=100$ & 982 & 319 & 2.727 \\
ours & edlm, $N=2000$, $\lambda=1000$ & 982 & 319 & 2.727 \\
D-CBG & $\gamma=10$, exact & 892 & 416 & 4.885 \\
D-CBG & $\gamma=8$, exact & 910 & 449 & 4.837 \\
\bottomrule
\end{tabular}}
\caption{UDLM: the Pareto-optimal settings of each method on (novel count, novel-ring count).}
\label{tab:frontier-ring_count-UDLM}
\end{table}

\section{Beyond the mean}\label{sec:beyond-mean-ring_count}

Everything above compares \emph{means}, which is the paper's metric but not the
one a screening pipeline faces. Nobody consumes the mean of a generated library;
they take the usable molecules out of it. And the mean is blind to the axis these
methods differ on most: from the same 1{,}024 generated sequences our runs yield
several times more novel molecules than D-CBG's.

\subsection{The primary comparison: usable molecules per fixed budget}

Both methods generated exactly 1{,}024 sequences. So simply count the novel
molecules that clear a quality bar. Quality and quantity both enter --- a run
that yields many poor novel molecules scores nothing, and one that yields a
handful of excellent ones scores little.

This number needs no correction of any kind. In particular it must \textbf{not}
be normalised by the novel count, nor size-matched against the competitor:
converting a fixed budget into more usable molecules \emph{is} the result, and
dividing it out would erase precisely what is being measured.


\subsubsection{The whole grid, on the primary metric}

The tables above pick a best setting per method, which is what a headline needs
but not what a reader should have to trust. Below is the raw grid: hits at the
headline bar at every $(N, \lambda)$, both modes, nothing averaged and nothing
filtered. It shows where in the grid the hits live --- and that the surface is
not monotone in either axis, so a single ``best'' cell is partly a draw from run
noise.

\begin{table}[htbp]
\centering
\scriptsize
\setlength{\tabcolsep}{3.5pt}
\resizebox{\ifdim\width>\linewidth\linewidth\else\width\fi}{!}{%
\begin{tabular}{rrrrrrrrrrr}
\toprule
$N$ & marg.\,0 & marg.\,5 & marg.\,10 & marg.\,30 & marg.\,1000 & edlm.\,0 & edlm.\,5 & edlm.\,10 & edlm.\,30 & edlm.\,1000 \\
\midrule
10 & 27 & 46 & 46 & 46 & 46 & 36 & 53 & 53 & 53 & 53 \\
30 & 28 & 54 & 56 & 56 & 56 & 43 & 53 & 51 & 51 & 51 \\
100 & 30 & 74 & 76 & 75 & 75 & 27 & 68 & 74 & 74 & 74 \\
300 & 33 & 85 & 92 & 93 & 93 & 25 & 82 & 80 & 80 & 80 \\
500 & 28 & 93 & 92 & 92 & 92 & 23 & 80 & 82 & 81 & 81 \\
1000 & 23 & 107 & 112 & 112 & 112 & 32 & 101 & 114 & 114 & 114 \\
2000 & 26 & 122 & 107 & 109 & 109 & 42 & 125 & 124 & 125 & 125 \\
\bottomrule
\end{tabular}}
\\[0.4em] \begin{minipage}{0.92\textwidth}\footnotesize Novel molecules clearing the bar, out of the 1{,}024 generated. Every cell is a single run, seed 1; nothing is averaged. $\lambda \in \{0.5, 1, 2, 100\}$ are omitted to fit the page. \textbf{Not} filtered by \texttt{MIN\_NOVEL}, unlike the Pareto table --- a collapsed run simply scores few hits here, which is the honest reading, whereas a mean over 2 molecules is not. \textbf{D-CBG's best on this bar is 76} (exact $\gamma=1$), for comparison.\end{minipage}
\caption{MDLM: hits at ring count $\geq 5$ across the whole $(N, \lambda)$ grid. Columns are (mixture mode) $\times$ $\lambda$.}
\label{tab:hits-grid-ring_count-MDLM}
\end{table}

\begin{table}[htbp]
\centering
\scriptsize
\setlength{\tabcolsep}{3.5pt}
\resizebox{\ifdim\width>\linewidth\linewidth\else\width\fi}{!}{%
\begin{tabular}{rrrrrrrrrrr}
\toprule
$N$ & marg.\,0 & marg.\,5 & marg.\,10 & marg.\,30 & marg.\,1000 & edlm.\,0 & edlm.\,5 & edlm.\,10 & edlm.\,30 & edlm.\,1000 \\
\midrule
10 & 4 & 16 & 15 & 15 & 15 & 4 & 8 & 7 & 7 & 7 \\
30 & 10 & 8 & 12 & 12 & 12 & 6 & 12 & 11 & 11 & 11 \\
100 & 5 & 18 & 16 & 16 & 16 & 5 & 7 & 9 & 9 & 9 \\
300 & 7 & 21 & 18 & 19 & 19 & 8 & 31 & 17 & 16 & 16 \\
500 & 4 & 25 & 26 & 28 & 28 & 5 & 19 & 25 & 24 & 24 \\
1000 & 8 & 29 & 29 & 31 & 31 & 10 & 38 & 30 & 30 & 30 \\
2000 & 2 & 29 & 33 & 34 & 34 & 7 & 49 & 54 & 54 & 54 \\
\bottomrule
\end{tabular}}
\\[0.4em] \begin{minipage}{0.92\textwidth}\footnotesize Novel molecules clearing the bar, out of the 1{,}024 generated. Every cell is a single run, seed 1; nothing is averaged. $\lambda \in \{0.5, 1, 2, 100\}$ are omitted to fit the page. \textbf{Not} filtered by \texttt{MIN\_NOVEL}, unlike the Pareto table --- a collapsed run simply scores few hits here, which is the honest reading, whereas a mean over 2 molecules is not. \textbf{D-CBG's best on this bar is 246} (exact $\gamma=8$), for comparison.\end{minipage}
\caption{UDLM: hits at ring count $\geq 5$ across the whole $(N, \lambda)$ grid. Columns are (mixture mode) $\times$ $\lambda$.}
\label{tab:hits-grid-ring_count-UDLM}
\end{table}

\begin{table}[htbp]
\centering
\footnotesize
\setlength{\tabcolsep}{4pt}
\resizebox{\ifdim\width>\linewidth\linewidth\else\width\fi}{!}{%
\begin{tabular}{lll rrrrr}
\toprule
model & method & setting & novel & mean & $\geq 4$ & $\geq 5$ & $\geq 6$ \\
\midrule
MDLM & \texttt{marginal} & $N=2000$, $\lambda=5$ & 354 & 3.715 & \textbf{177} (1.3$\times$) & \textbf{122} (1.6$\times$) & \textbf{52} (1.5$\times$) \\
MDLM & \texttt{edlm} & $N=2000$, $\lambda=5$ & 367 & 3.793 & \textbf{184} (1.3$\times$) & \textbf{125} (1.6$\times$) & \textbf{66} (1.9$\times$) \\
MDLM & D-CBG & $\gamma=1$, exact & 149 & 4.745 & 137 & 76 & 34 \\
UDLM & \texttt{marginal} & $N=2000$, $\lambda=30$ & 277 & 2.581 & 59 (0.1$\times$) & 34 (0.1$\times$) & 8 (0.1$\times$) \\
UDLM & \texttt{edlm} & $N=2000$, $\lambda=10$ & 311 & 2.717 & 82 (0.2$\times$) & 54 (0.2$\times$) & 7 (0.1$\times$) \\
UDLM & D-CBG & $\gamma=8$, exact & 449 & 4.837 & 444 & 246 & 93 \\
\bottomrule
\end{tabular}}
\\[0.4em] \begin{minipage}{0.92\textwidth}\footnotesize No size matching and no normalisation --- every run spent the same generation budget, so these are directly comparable counts.\end{minipage}
\caption{Novel molecules clearing each ring count bar, out of the 1{,}024 sequences each run generated. The two mixture-sampling modes are kept apart rather than merged into a single \emph{ours} row, since they are different samplers. Each is at its own best setting, ranked on the middle bar; the multiplier is against D-CBG on the same model.}
\label{tab:hits-ring_count}
\end{table}


\subsection{Secondary: shape comparisons, which discard the counts}

Two standard distribution-free comparisons. Both deliberately ignore how many
molecules each method produced, so neither can see the advantage above; they are
here to answer a different question --- whether our \emph{distribution} is
better, or only our yield.

\begin{itemize}
  \item \textbf{Rank against rank.} Sort both sets and compare like for like:
        best against best, tenth percentile against tenth. The fraction of
        quantiles where we lead is the \emph{quantile win rate}; 1.0 would be
        \textbf{first-order stochastic dominance}. Comparing the $k$-th
        \emph{rank} directly would need equal sample sizes, so this compares at
        matched quantiles, the size-independent form of the same idea.
  \item \textbf{Draw against draw.} $A_{12}$, the Vargha--Delaney probability of
        superiority: one molecule drawn at random from each, how often is ours
        better? The normalised Mann--Whitney $U$; $0.5$ is indistinguishable.
\end{itemize}

The \emph{matched} column subsamples the larger novel set down to the smaller.
It is a \textbf{decomposition, not a fairness correction} --- it separates
``we win because our tail is better'' from ``we win because our pool is deeper''.
The counterfactual it describes never occurs in practice, so it should not be
read as the fair comparison; the table above is.

\begin{table}[htbp]
\centering
\small
\resizebox{\ifdim\width>\linewidth\linewidth\else\width\fi}{!}{%
\begin{tabular}{ll rrrrr rr}
\toprule
model & method & novel & mean & max & top10 & top10 matched & q-win & $A_{12}$ \\
\midrule
MDLM & ours & 304 & 3.615 & 12.000 & 10.100 & 8.814 & 0.43 & 0.39 \\
MDLM & D-CBG & 111 & 3.856 & 10.000 & 7.600 & 7.600 & -- & -- \\
UDLM & ours & 313 & 2.540 & 10.000 & 6.800 & 6.800 & 0.00 & 0.08 \\
UDLM & D-CBG & 416 & 4.885 & 10.000 & 9.100 & 8.575 & -- & -- \\
\bottomrule
\end{tabular}}
\\[0.4em] \begin{minipage}{0.92\textwidth}\footnotesize \emph{max} and \emph{top-K} are at the fixed 1,024-sequence budget and so include our deeper novel pool. \emph{matched} removes that, and answers only ``quality or quantity?.\end{minipage}
\caption{Each method at its own best setting by top-10. \emph{q-win} and $A_{12}$ sit on the ours row and describe ours against D-CBG.}
\label{tab:beyond-ring_count}
\end{table}

\section{Validity}\label{sec:validity-ring_count}

Validity is the axis on which the two mechanisms differ most sharply, and it is
what drives the curve directions of the previous section, so it is worth
reporting on its own.

Our reward sees a \emph{clean} $x_0$ candidate and returns
\texttt{invalid\_reward} $=0$ when RDKit cannot parse it. Raising $\lambda$
therefore pushes probability towards parseable molecules before it pushes
towards high ring count: an invalid candidate is treated as a ring count-zero one and
suppressed. D-CBG instead perturbs the per-position logits with the gradient of
a noisy classifier, which carries no such term --- nothing in it prefers a
parseable string.

\begin{table}[htbp]
\centering
\small
\resizebox{\ifdim\width>\linewidth\linewidth\else\width\fi}{!}{%
\begin{tabular}{rrrrrrrrrr}
\toprule
$N$ & $\lambda=0$ & $\lambda=0.5$ & $\lambda=1$ & $\lambda=2$ & $\lambda=5$ & $\lambda=10$ & $\lambda=30$ & $\lambda=100$ & $\lambda=1000$ \\
\midrule
10 & 547 & 580 & 600 & 610 & 628 & 629 & 629 & 629 & 629 \\
30 & 542 & 569 & 595 & 623 & 637 & 637 & 637 & 637 & 637 \\
100 & 532 & 568 & 600 & 638 & 659 & 657 & 658 & 658 & 658 \\
300 & 552 & 580 & 600 & 657 & 683 & 700 & 700 & 700 & 700 \\
500 & 558 & 584 & 618 & 650 & 684 & 696 & 695 & 695 & 695 \\
1000 & 549 & 583 & 613 & 653 & 697 & 701 & 702 & 702 & 702 \\
2000 & 544 & 565 & 611 & 654 & 716 & 731 & 734 & 734 & 734 \\
\bottomrule
\end{tabular}}
\caption{MDLM, \texttt{marginal}: valid molecules out of the 1{,}024 generated.}
\label{tab:valid-ring_count-MDLM-marginal}
\end{table}

\begin{table}[htbp]
\centering
\small
\resizebox{\ifdim\width>\linewidth\linewidth\else\width\fi}{!}{%
\begin{tabular}{rrrrrrrrrr}
\toprule
$N$ & $\lambda=0$ & $\lambda=0.5$ & $\lambda=1$ & $\lambda=2$ & $\lambda=5$ & $\lambda=10$ & $\lambda=30$ & $\lambda=100$ & $\lambda=1000$ \\
\midrule
10 & 564 & 590 & 605 & 622 & 629 & 629 & 629 & 629 & 629 \\
30 & 558 & 590 & 608 & 622 & 639 & 640 & 640 & 640 & 640 \\
100 & 553 & 598 & 631 & 656 & 704 & 699 & 698 & 698 & 698 \\
300 & 556 & 594 & 620 & 659 & 680 & 696 & 696 & 696 & 696 \\
500 & 549 & 581 & 614 & 658 & 696 & 704 & 704 & 704 & 704 \\
1000 & 542 & 577 & 631 & 678 & 717 & 728 & 730 & 730 & 730 \\
2000 & 559 & 591 & 633 & 694 & 766 & 754 & 754 & 754 & 754 \\
\bottomrule
\end{tabular}}
\caption{MDLM, \texttt{edlm}: valid molecules out of the 1{,}024 generated.}
\label{tab:valid-ring_count-MDLM-edlm}
\end{table}

\begin{table}[htbp]
\centering
\small
\resizebox{\ifdim\width>\linewidth\linewidth\else\width\fi}{!}{%
\begin{tabular}{rrrrrrrrrr}
\toprule
$N$ & $\lambda=0$ & $\lambda=0.5$ & $\lambda=1$ & $\lambda=2$ & $\lambda=5$ & $\lambda=10$ & $\lambda=30$ & $\lambda=100$ & $\lambda=1000$ \\
\midrule
10 & 682 & 731 & 770 & 832 & 858 & 866 & 867 & 867 & 867 \\
30 & 695 & 759 & 791 & 867 & 902 & 922 & 922 & 922 & 922 \\
100 & 686 & 743 & 798 & 863 & 916 & 921 & 921 & 921 & 921 \\
300 & 672 & 733 & 774 & 872 & 940 & 957 & 961 & 961 & 961 \\
500 & 691 & 744 & 792 & 862 & 943 & 953 & 951 & 951 & 951 \\
1000 & 695 & 748 & 792 & 880 & 938 & 954 & 961 & 961 & 961 \\
2000 & 694 & 763 & 817 & 873 & 954 & 972 & 974 & 974 & 974 \\
\bottomrule
\end{tabular}}
\caption{UDLM, \texttt{marginal}: valid molecules out of the 1{,}024 generated.}
\label{tab:valid-ring_count-UDLM-marginal}
\end{table}

\begin{table}[htbp]
\centering
\small
\resizebox{\ifdim\width>\linewidth\linewidth\else\width\fi}{!}{%
\begin{tabular}{rrrrrrrrrr}
\toprule
$N$ & $\lambda=0$ & $\lambda=0.5$ & $\lambda=1$ & $\lambda=2$ & $\lambda=5$ & $\lambda=10$ & $\lambda=30$ & $\lambda=100$ & $\lambda=1000$ \\
\midrule
10 & 657 & 731 & 786 & 845 & 889 & 890 & 890 & 890 & 890 \\
30 & 675 & 744 & 802 & 879 & 912 & 919 & 919 & 919 & 919 \\
100 & 679 & 746 & 813 & 879 & 944 & 953 & 953 & 953 & 953 \\
300 & 693 & 765 & 810 & 901 & 969 & 978 & 979 & 979 & 979 \\
500 & 686 & 775 & 806 & 882 & 966 & 979 & 980 & 980 & 980 \\
1000 & 685 & 748 & 808 & 901 & 978 & 984 & 986 & 986 & 986 \\
2000 & 685 & 766 & 824 & 902 & 981 & 979 & 982 & 982 & 982 \\
\bottomrule
\end{tabular}}
\caption{UDLM, \texttt{edlm}: valid molecules out of the 1{,}024 generated.}
\label{tab:valid-ring_count-UDLM-edlm}
\end{table}

\begin{table}[htbp]
\centering
\small
\resizebox{\ifdim\width>\linewidth\linewidth\else\width\fi}{!}{%
\begin{tabular}{ll rrr l}
\toprule
model & method & weakest dial & best & strongest dial & trend \\
\midrule
MDLM & ours ($N=2000$) & 544 & 734 & 734 & \textbf{rises} \\
MDLM & D-CBG exact & 188 & 188 & 2 & falls \\
MDLM & D-CBG approx. & 493 & 493 & 88 & falls \\
UDLM & ours ($N=2000$) & 694 & 974 & 974 & \textbf{rises} \\
UDLM & D-CBG exact & 807 & 910 & 892 & rises \\
UDLM & D-CBG approx. & 992 & 992 & 792 & falls \\
\bottomrule
\end{tabular}}
\\[0.4em] \begin{minipage}{0.92\textwidth}\footnotesize This is the mechanism behind the opposite curve directions: our dial buys validity on the way to buying property, so the novel count rises with it, while D-CBG spends validity and the novel count falls. On MDLM the D-CBG exact arm reaches zero valid molecules at $\gamma \geq 6$, which is why the paper does not report it there.\end{minipage}
\caption{Valid molecules out of 1{,}024 as the strength dial is turned up. Ours is read along $\lambda$ at the largest $N$; D-CBG along $\gamma$.}
\label{tab:valid-vs-dial-ring_count}
\end{table}

\section{The shape of each curve}

Plot novel count on the $x$-axis against novel-ring count mean on the $y$-axis and
sweep each method's strength dial --- $\lambda$ for ours, $\gamma$ for D-CBG.
The two methods trace curves in \emph{opposite directions}, and that is a
property worth reporting on its own, independently of which sits higher.

\begin{itemize}
  \item \textbf{D-CBG runs up and to the left.} Raising $\gamma$ buys
        property by spending novelty: the dial is a genuine trade-off the user
        has to tune, and the original paper tunes it per model and per property.
  \item \textbf{Ours runs up and to the right.} Raising $\lambda$ improves
        \emph{both} axes at once. The reason is visible in the validity
        column: the first thing $\lambda$ buys is valid molecules, which raises
        the novel count, and only then does it buy property. There is no
        trade-off to tune --- the only cost is compute.
\end{itemize}

The correlation between the two axes along each sweep reduces this to one
number. Crucially it is measured at \emph{fixed} $N$, so the up-right shape
is not an artifact of spending more Monte Carlo budget: along every row below,
only the strength dial moves.

\begin{table}[htbp]
\centering
\small
\resizebox{\ifdim\width>\linewidth\linewidth\else\width\fi}{!}{%
\begin{tabular}{ll r rr rl}
\toprule
model & method & pts & weakest dial & strongest dial & corr & shape \\
\midrule
MDLM & ours, $N=100$ & 9 & (211, 2.346) & (279, 3.108) & +0.98 & \textbf{up-right} \\
MDLM & ours, $N=300$ & 9 & (220, 2.468) & (316, 3.367) & +0.93 & \textbf{up-right} \\
MDLM & ours, $N=1000$ & 9 & (204, 2.426) & (369, 3.493) & +0.96 & \textbf{up-right} \\
MDLM & ours, $N=2000$ & 9 & (190, 2.379) & (362, 3.627) & +0.96 & \textbf{up-right} \\
MDLM & D-CBG approx. & 6 & (209, 2.809) & (67, 4.597) & -0.99 & up-left \\
UDLM & ours, $N=100$ & 9 & (194, 1.660) & (265, 1.932) & +0.95 & \textbf{up-right} \\
UDLM & ours, $N=300$ & 9 & (193, 1.637) & (288, 2.125) & +0.96 & \textbf{up-right} \\
UDLM & ours, $N=1000$ & 9 & (186, 1.543) & (308, 2.351) & +0.97 & \textbf{up-right} \\
UDLM & ours, $N=2000$ & 9 & (207, 1.469) & (277, 2.581) & +0.96 & \textbf{up-right} \\
UDLM & D-CBG exact & 6 & (275, 4.033) & (416, 4.885) & +0.86 & \textbf{up-right} \\
UDLM & D-CBG approx. & 6 & (154, 2.519) & (355, 4.780) & +0.87 & \textbf{up-right} \\
\bottomrule
\end{tabular}}
\\[0.4em] \begin{minipage}{0.92\textwidth}\footnotesize Measured at fixed $N$, so the up-right shape is not an effect of raising the Monte Carlo budget --- only the strength dial moves along each row.\end{minipage}
\caption{Direction of each curve in (novel count, novel-ring count) space. \emph{corr} is between the two axes along the sweep: positive means the dial improves both, negative means it trades one for the other.}
\label{tab:shape-ring_count}
\end{table}

\begin{table}[htbp]
\centering
\scriptsize
\setlength{\tabcolsep}{3.5pt}
\resizebox{\ifdim\width>\linewidth\linewidth\else\width\fi}{!}{%
\begin{tabular}{lccccccccc}
\toprule
ours & $\lambda=0$ & $\lambda=0.5$ & $\lambda=1$ & $\lambda=2$ & $\lambda=5$ & $\lambda=10$ & $\lambda=30$ & $\lambda=100$ & $\lambda=1000$ \\
\midrule
$N=10$ & (216, 2.477) & (233, 2.751) & (244, 2.824) & (243, 2.815) & (254, 2.819) & (256, 2.820) & (256, 2.820) & (256, 2.820) & (256, 2.820) \\
$N=30$ & (211, 2.365) & (232, 2.552) & (245, 2.657) & (265, 2.683) & (276, 2.681) & (288, 2.639) & (289, 2.633) & (289, 2.633) & (289, 2.633) \\
$N=100$ & (211, 2.346) & (232, 2.573) & (259, 2.822) & (281, 2.979) & (287, 3.080) & (279, 3.129) & (279, 3.108) & (279, 3.108) & (279, 3.108) \\
$N=300$ & (220, 2.468) & (236, 2.716) & (241, 3.033) & (281, 3.260) & (293, 3.345) & (316, 3.351) & (316, 3.367) & (316, 3.367) & (316, 3.367) \\
$N=500$ & (228, 2.548) & (245, 2.698) & (272, 2.952) & (294, 3.173) & (331, 3.311) & (334, 3.413) & (334, 3.410) & (334, 3.410) & (334, 3.410) \\
$N=1000$ & (204, 2.426) & (229, 2.603) & (257, 2.911) & (281, 3.310) & (361, 3.452) & (367, 3.499) & (369, 3.493) & (369, 3.493) & (369, 3.493) \\
$N=2000$ & (190, 2.379) & (204, 2.672) & (245, 3.045) & (304, 3.615) & (354, 3.715) & (357, 3.611) & (362, 3.627) & (362, 3.627) & (362, 3.627) \\
\bottomrule
\end{tabular}}
\caption{MDLM: the (novel, novel-ring count) point at each $\lambda$, for every $N$. Reading a row left to right traces that curve.}
\label{tab:curve-ring_count-MDLM}
\end{table}

\begin{table}[htbp]
\centering
\scriptsize
\setlength{\tabcolsep}{3.5pt}
\resizebox{\ifdim\width>\linewidth\linewidth\else\width\fi}{!}{%
\begin{tabular}{lllllll}
\toprule
D-CBG &  &  &  &  &  &  \\
\midrule
exact & $\gamma=1$: (149, 4.745) &  &  &  &  &  \\
approx. & $\gamma=1$: (209, 2.809) & $\gamma=2$: (176, 2.966) & $\gamma=3$: (150, 3.400) & $\gamma=5$: (111, 3.856) & $\gamma=8$: (86, 4.244) & $\gamma=10$: (67, 4.597) \\
\bottomrule
\end{tabular}}
\caption{MDLM: the same for D-CBG, sweeping $\gamma$. The novel count falls as the property rises --- the opposite direction.}
\label{tab:curve-cbg-ring_count-MDLM}
\end{table}

\begin{table}[htbp]
\centering
\scriptsize
\setlength{\tabcolsep}{3.5pt}
\resizebox{\ifdim\width>\linewidth\linewidth\else\width\fi}{!}{%
\begin{tabular}{lccccccccc}
\toprule
ours & $\lambda=0$ & $\lambda=0.5$ & $\lambda=1$ & $\lambda=2$ & $\lambda=5$ & $\lambda=10$ & $\lambda=30$ & $\lambda=100$ & $\lambda=1000$ \\
\midrule
$N=10$ & (189, 1.656) & (208, 1.865) & (230, 1.935) & (251, 1.873) & (243, 1.868) & (241, 1.884) & (241, 1.884) & (241, 1.884) & (241, 1.884) \\
$N=30$ & (195, 1.708) & (214, 1.785) & (202, 1.787) & (242, 1.884) & (251, 1.789) & (264, 1.799) & (265, 1.804) & (265, 1.804) & (265, 1.804) \\
$N=100$ & (194, 1.660) & (214, 1.776) & (229, 1.825) & (240, 1.846) & (286, 1.920) & (267, 1.936) & (265, 1.932) & (265, 1.932) & (265, 1.932) \\
$N=300$ & (193, 1.637) & (207, 1.768) & (212, 1.764) & (240, 1.954) & (268, 2.168) & (292, 2.110) & (288, 2.125) & (288, 2.125) & (288, 2.125) \\
$N=500$ & (196, 1.592) & (217, 1.627) & (222, 1.721) & (249, 2.008) & (278, 2.317) & (284, 2.162) & (279, 2.211) & (279, 2.211) & (279, 2.211) \\
$N=1000$ & (186, 1.543) & (206, 1.646) & (228, 1.842) & (241, 2.129) & (288, 2.302) & (294, 2.364) & (308, 2.351) & (308, 2.351) & (308, 2.351) \\
$N=2000$ & (207, 1.469) & (234, 1.654) & (243, 1.733) & (256, 1.953) & (275, 2.524) & (278, 2.572) & (277, 2.581) & (277, 2.581) & (277, 2.581) \\
\bottomrule
\end{tabular}}
\caption{UDLM: the (novel, novel-ring count) point at each $\lambda$, for every $N$. Reading a row left to right traces that curve.}
\label{tab:curve-ring_count-UDLM}
\end{table}

\begin{table}[htbp]
\centering
\scriptsize
\setlength{\tabcolsep}{3.5pt}
\resizebox{\ifdim\width>\linewidth\linewidth\else\width\fi}{!}{%
\begin{tabular}{lllllll}
\toprule
D-CBG &  &  &  &  &  &  \\
\midrule
exact & $\gamma=1$: (275, 4.033) & $\gamma=2$: (307, 4.619) & $\gamma=3$: (345, 4.557) & $\gamma=5$: (400, 4.695) & $\gamma=8$: (449, 4.837) & $\gamma=10$: (416, 4.885) \\
approx. & $\gamma=1$: (154, 2.519) & $\gamma=2$: (147, 3.184) & $\gamma=3$: (197, 3.893) & $\gamma=5$: (223, 4.435) & $\gamma=8$: (298, 4.661) & $\gamma=10$: (355, 4.780) \\
\bottomrule
\end{tabular}}
\caption{UDLM: the same for D-CBG, sweeping $\gamma$. The novel count falls as the property rises --- the opposite direction.}
\label{tab:curve-cbg-ring_count-UDLM}
\end{table}

\section{What the sweep says}
\textbf{MDLM.} Peak novel-ring count: ours 3.793 against D-CBG 4.745. Peak novel count: ours 379 against D-CBG 209. Each method wins one axis, so the comparison is a genuine frontier.
\textbf{UDLM.} Peak novel-ring count: ours 2.727 against D-CBG 4.885. Peak novel count: ours 319 against D-CBG 449. \emph{D-CBG dominates on both axes} --- there is no trade-off to report here, we simply lose.

\subsection{Why ring count behaves differently from QED}

The method can only \emph{select}, never \emph{invent}. Weights are
$\mathrm{softmax}(\lambda r)$ over $N$ candidates drawn from $x_\theta$, so the
best reachable reward at a step is the best value \emph{present among those $N$
candidates}. No $\lambda$ can produce a 5-ring molecule if none of the $N$
proposals has one.

That ceiling binds hard on ring count and barely at all on QED, for two reasons
that reinforce each other:

\begin{itemize}
  \item \textbf{The reward is integer valued}, so candidates tie and the tilt
        saturates early. Raising $\lambda$ past $\approx5$ changes nothing:
        the softmax is already a hard argmax over the tied top set.
  \item \textbf{High-ring molecules are rare under the base model} (QM9's mean
        ring count is 1.74), so the $N$ proposals rarely contain one.
\end{itemize}

D-CBG has no such ceiling: it perturbs the per-position logits with a classifier
gradient and can steer toward regions the base model would not have proposed.
That is exactly what the numbers show.

\textbf{The binding constraint is $N$, not $\lambda$.} Along the saturated
$\lambda$ ridge the property still climbs monotonically with $N$ and has not
flattened at $N=500$. Pushing $N$ further is the one lever with headroom left;
raising $\lambda$ is spent.

\end{document}